% --- Archon physics formalization source begin ---
% archon:physics
% archon:covers PhyXMiniProblems/problem_phyx_mini_0583.lean
% archon:source-report reports/phyx_mini/problem_phyx_mini_0583.source.json

\chapter{Physics problem phyx\_mini\_0583}
\label{ch:PhyXMiniProblems_problem_phyx_mini_0583}

\paragraph{Problem source.}
\#\# Physical scenario

The two-dimensional, infinite corral of figure is square, with edge length \$L = 150 \textbackslash{}, \textbackslash{}text\{pm\}\$. A square probe is centered at \$xy\$ coordinates \$(0.200L, 0.800L)\$ and has an \$x\$ width of \$5.00 \textbackslash{}, \textbackslash{}text\{pm\}\$ and a \$y\$ width of \$5.00 \textbackslash{}, \textbackslash{}text\{pm\}\$.

\#\# Question

What is the probability of detection if the electron is in the \$E\_\{1,3\}\$ energy state?

\#\# Answer choices

- A: A: 1.4 \textbackslash{}times 10\textasciicircum{}\{-2\} \textbackslash{}

- B: B: 2.8 \textbackslash{}times 10\textasciicircum{}\{-3\} \textbackslash{}

- C: C: 1.4 \textbackslash{}times 10\textasciicircum{}\{-3\} \textbackslash{}

- D: D: 7.0 \textbackslash{}times 10\textasciicircum{}\{-4\}

\#\# Figure caption (auxiliary; use the image as primary evidence)

The image is a diagram that features a square labeled as "Probe" towards the top left inside a larger square with a thick blue border. The probe is represented by a smaller square and is marked with a line pointing to the text "Probe." The axes are labeled \textbackslash{}(x\textbackslash{}) (horizontal) and \textbackslash{}(y\textbackslash{}) (vertical), suggesting a coordinate plane. The overall layout suggests a spatial or positional context for the "Probe" within the larger square boundary.

\#\# Dataset metadata

Quantum Phenomena; Spatial Relation Reasoning, Numerical Reasoning

\paragraph{Recorded answer/context.}
C: C: 1.4 \textbackslash{}times 10\textasciicircum{}\{-3\} \textbackslash{}

\paragraph{Figure/image path.}
/root/proposal\_for\_physic/science-mango/phyx\_mini\_run/phyx\_data/test\_image/583.png

\paragraph{Formalization target.}
create a compiling Lean file with sorry bodies at `PhyXMiniProblems/problem\_phyx\_mini\_0583.lean`. The Lean declarations must preserve the physical quantities, dimensions or dimensional roles, figure labels, governing-law hypotheses, and final relation expressed by this problem.
Use Mathlib/Physlib names found through LeanExplore where available. If a physics API is missing, introduce faithful local abstractions rather than scalar placeholder aliases.

\begin{theorem}[Physics formalization target]
\label{thm:physics:phyx_mini_0583:target}
\lean{PhyXMiniProblems.ProblemPhyXMini0583.problem_phyx_mini_0583}
\uses{def:physics:phyx-mini-0583:phyxminiproblems-problemphyxmini0583-squarecorralexperiment, def:physics:phyx-mini-0583:phyxminiproblems-problemphyxmini0583-satisfiesinfinitesquarecorralphysics, def:physics:phyx-mini-0583:phyxminiproblems-problemphyxmini0583-matchesproblemstatementandfigure, def:physics:phyx-mini-0583:phyxminiproblems-problemphyxmini0583-matchesanswerchoice}
The assigned autoformalize agent should translate this physics problem into Lean declarations in the covered file, with theorem and lemma proof bodies written as `by sorry`.
\end{theorem}
\begin{proof}
This is an autoformalization task, not a proof task. Produce statements that can later be proved by the physics prover without weakening the physical model.
\end{proof}

% --- Archon declaration topology begin ---
\section{Declaration topology}
\begin{definition}[Length Magnitude]
  \label{def:physics:phyx-mini-0583:phyxminiproblems-problemphyxmini0583-lengthmagnitude}
  \lean{PhyXMiniProblems.ProblemPhyXMini0583.LengthMagnitude}
  A nonnegative physical length, independent of the chosen unit system
\end{definition}
\begin{proof}
  This is the declared model definition of Length Magnitude.
\end{proof}
\begin{definition}[Signed Length Quantity]
  \label{def:physics:phyx-mini-0583:phyxminiproblems-problemphyxmini0583-signedlengthquantity}
  \lean{PhyXMiniProblems.ProblemPhyXMini0583.SignedLengthQuantity}
  A signed physical position along one Cartesian axis
\end{definition}
\begin{proof}
  This is the declared model definition of Signed Length Quantity.
\end{proof}
\begin{definition}[length Readout]
  \label{def:physics:phyx-mini-0583:phyxminiproblems-problemphyxmini0583-lengthreadout}
  \lean{PhyXMiniProblems.ProblemPhyXMini0583.lengthReadout}
  \uses{def:physics:phyx-mini-0583:phyxminiproblems-problemphyxmini0583-lengthmagnitude}
  Read a nonnegative physical length in a selected Physlib length unit
\end{definition}
\begin{proof}
  This is the declared model definition of length Readout.
\end{proof}
\begin{definition}[signed Length Readout]
  \label{def:physics:phyx-mini-0583:phyxminiproblems-problemphyxmini0583-signedlengthreadout}
  \lean{PhyXMiniProblems.ProblemPhyXMini0583.signedLengthReadout}
  \uses{def:physics:phyx-mini-0583:phyxminiproblems-problemphyxmini0583-signedlengthquantity}
  Read a signed Cartesian position in a selected Physlib length unit
\end{definition}
\begin{proof}
  This is the declared model definition of signed Length Readout.
\end{proof}
\begin{definition}[length In Picometers]
  \label{def:physics:phyx-mini-0583:phyxminiproblems-problemphyxmini0583-lengthinpicometers}
  \lean{PhyXMiniProblems.ProblemPhyXMini0583.lengthInPicometers}
  \uses{def:physics:phyx-mini-0583:phyxminiproblems-problemphyxmini0583-lengthmagnitude, def:physics:phyx-mini-0583:phyxminiproblems-problemphyxmini0583-lengthreadout}
  Picometre readout of a physical length
\end{definition}
\begin{proof}
  This is the declared model definition of length In Picometers.
\end{proof}
\begin{definition}[signed Length In Picometers]
  \label{def:physics:phyx-mini-0583:phyxminiproblems-problemphyxmini0583-signedlengthinpicometers}
  \lean{PhyXMiniProblems.ProblemPhyXMini0583.signedLengthInPicometers}
  \uses{def:physics:phyx-mini-0583:phyxminiproblems-problemphyxmini0583-signedlengthquantity, def:physics:phyx-mini-0583:phyxminiproblems-problemphyxmini0583-signedlengthreadout}
  Picometre readout of a signed Cartesian position
\end{definition}
\begin{proof}
  This is the declared model definition of signed Length In Picometers.
\end{proof}
\begin{definition}[Quantum Particle Kind]
  \label{def:physics:phyx-mini-0583:phyxminiproblems-problemphyxmini0583-quantumparticlekind}
  \lean{PhyXMiniProblems.ProblemPhyXMini0583.QuantumParticleKind}
  The particle species named in the problem
\end{definition}
\begin{proof}
  This is the declared model definition of Quantum Particle Kind.
\end{proof}
\begin{definition}[Corral Confinement]
  \label{def:physics:phyx-mini-0583:phyxminiproblems-problemphyxmini0583-corralconfinement}
  \lean{PhyXMiniProblems.ProblemPhyXMini0583.CorralConfinement}
  The idealized confinement represented by the blue square boundary
\end{definition}
\begin{proof}
  This is the declared model definition of Corral Confinement.
\end{proof}
\begin{definition}[Cartesian Axis Label]
  \label{def:physics:phyx-mini-0583:phyxminiproblems-problemphyxmini0583-cartesianaxislabel}
  \lean{PhyXMiniProblems.ProblemPhyXMini0583.CartesianAxisLabel}
  The two Cartesian labels printed on the image axes
\end{definition}
\begin{proof}
  This is the declared model definition of Cartesian Axis Label.
\end{proof}
\begin{definition}[Figure Shape]
  \label{def:physics:phyx-mini-0583:phyxminiproblems-problemphyxmini0583-figureshape}
  \lean{PhyXMiniProblems.ProblemPhyXMini0583.FigureShape}
  The geometric shape of the confining boundary and of the probe
\end{definition}
\begin{proof}
  This is the declared model definition of Figure Shape.
\end{proof}
\begin{definition}[Infinite Square Corral Mode]
  \label{def:physics:phyx-mini-0583:phyxminiproblems-problemphyxmini0583-infinitesquarecorralmode}
  \lean{PhyXMiniProblems.ProblemPhyXMini0583.InfiniteSquareCorralMode}
  Quantum numbers in the x and y directions for a square-corral mode
\end{definition}
\begin{proof}
  This is the declared model definition of Infinite Square Corral Mode.
\end{proof}
\begin{definition}[Square Position Probe]
  \label{def:physics:phyx-mini-0583:phyxminiproblems-problemphyxmini0583-squarepositionprobe}
  \lean{PhyXMiniProblems.ProblemPhyXMini0583.SquarePositionProbe}
  \uses{def:physics:phyx-mini-0583:phyxminiproblems-problemphyxmini0583-lengthmagnitude, def:physics:phyx-mini-0583:phyxminiproblems-problemphyxmini0583-signedlengthquantity}
  A physical square position probe with a center and full widths
\end{definition}
\begin{proof}
  This is the declared model definition of Square Position Probe.
\end{proof}
\begin{definition}[Square Corral Figure]
  \label{def:physics:phyx-mini-0583:phyxminiproblems-problemphyxmini0583-squarecorralfigure}
  \lean{PhyXMiniProblems.ProblemPhyXMini0583.SquareCorralFigure}
  \uses{def:physics:phyx-mini-0583:phyxminiproblems-problemphyxmini0583-cartesianaxislabel, def:physics:phyx-mini-0583:phyxminiproblems-problemphyxmini0583-figureshape}
  Qualitative labels and shapes visible in the supplied image
\end{definition}
\begin{proof}
  This is the declared model definition of Square Corral Figure.
\end{proof}
\begin{definition}[Square Corral Experiment]
  \label{def:physics:phyx-mini-0583:phyxminiproblems-problemphyxmini0583-squarecorralexperiment}
  \lean{PhyXMiniProblems.ProblemPhyXMini0583.SquareCorralExperiment}
  \uses{def:physics:phyx-mini-0583:phyxminiproblems-problemphyxmini0583-lengthmagnitude, def:physics:phyx-mini-0583:phyxminiproblems-problemphyxmini0583-quantumparticlekind, def:physics:phyx-mini-0583:phyxminiproblems-problemphyxmini0583-corralconfinement, def:physics:phyx-mini-0583:phyxminiproblems-problemphyxmini0583-infinitesquarecorralmode, def:physics:phyx-mini-0583:phyxminiproblems-problemphyxmini0583-squarepositionprobe, def:physics:phyx-mini-0583:phyxminiproblems-problemphyxmini0583-squarecorralfigure}
  The electron experiment, including a representative wavefunction on numerical picometre coordinates and its genuine Physlib two-dimensional Hilbert-space state. The dependent fields ensure that the representative is square integrable and denotes stateVector; the physical probability observable is a separate normalized measure until Born's law is imposed below
\end{definition}
\begin{proof}
  This is the declared model definition of Square Corral Experiment.
\end{proof}
\begin{definition}[x Picometers]
  \label{def:physics:phyx-mini-0583:phyxminiproblems-problemphyxmini0583-xpicometers}
  \lean{PhyXMiniProblems.ProblemPhyXMini0583.xPicometers}
  Numerical x coordinate in picometres
\end{definition}
\begin{proof}
  This is the declared model definition of x Picometers.
\end{proof}
\begin{definition}[y Picometers]
  \label{def:physics:phyx-mini-0583:phyxminiproblems-problemphyxmini0583-ypicometers}
  \lean{PhyXMiniProblems.ProblemPhyXMini0583.yPicometers}
  Numerical y coordinate in picometres
\end{definition}
\begin{proof}
  This is the declared model definition of y Picometers.
\end{proof}
\begin{definition}[square Corral Region Picometers]
  \label{def:physics:phyx-mini-0583:phyxminiproblems-problemphyxmini0583-squarecorralregionpicometers}
  \lean{PhyXMiniProblems.ProblemPhyXMini0583.squareCorralRegionPicometers}
  \uses{def:physics:phyx-mini-0583:phyxminiproblems-problemphyxmini0583-xpicometers, def:physics:phyx-mini-0583:phyxminiproblems-problemphyxmini0583-ypicometers}
  The closed square occupied by a corral whose edge readout is L
\end{definition}
\begin{proof}
  This is the declared model definition of square Corral Region Picometers.
\end{proof}
\begin{definition}[square Probe Region Picometers]
  \label{def:physics:phyx-mini-0583:phyxminiproblems-problemphyxmini0583-squareproberegionpicometers}
  \lean{PhyXMiniProblems.ProblemPhyXMini0583.squareProbeRegionPicometers}
  \uses{def:physics:phyx-mini-0583:phyxminiproblems-problemphyxmini0583-lengthinpicometers, def:physics:phyx-mini-0583:phyxminiproblems-problemphyxmini0583-signedlengthinpicometers, def:physics:phyx-mini-0583:phyxminiproblems-problemphyxmini0583-squarepositionprobe, def:physics:phyx-mini-0583:phyxminiproblems-problemphyxmini0583-xpicometers, def:physics:phyx-mini-0583:phyxminiproblems-problemphyxmini0583-ypicometers}
  The closed rectangular region covered by the physical square probe
\end{definition}
\begin{proof}
  This is the declared model definition of square Probe Region Picometers.
\end{proof}
\begin{definition}[infinite Square Corral Mode Amplitude Per Picometer]
  \label{def:physics:phyx-mini-0583:phyxminiproblems-problemphyxmini0583-infinitesquarecorralmodeamplitudeperpicometer}
  \lean{PhyXMiniProblems.ProblemPhyXMini0583.infiniteSquareCorralModeAmplitudePerPicometer}
  \uses{def:physics:phyx-mini-0583:phyxminiproblems-problemphyxmini0583-infinitesquarecorralmode, def:physics:phyx-mini-0583:phyxminiproblems-problemphyxmini0583-xpicometers, def:physics:phyx-mini-0583:phyxminiproblems-problemphyxmini0583-ypicometers, def:physics:phyx-mini-0583:phyxminiproblems-problemphyxmini0583-squarecorralregionpicometers}
  The normalized stationary wave-amplitude representative for an infinite square corral in numerical picometre coordinates. Inside the box it is (2/L) sin(n\_x π x/L) sin(n\_y π y/L) and it vanishes outside. This is a general governing formula; it contains no requested probability or answer
\end{definition}
\begin{proof}
  This is the declared model definition of infinite Square Corral Mode Amplitude Per Picometer.
\end{proof}
\begin{definition}[Satisfies Infinite Square Corral Physics]
  \label{def:physics:phyx-mini-0583:phyxminiproblems-problemphyxmini0583-satisfiesinfinitesquarecorralphysics}
  \lean{PhyXMiniProblems.ProblemPhyXMini0583.SatisfiesInfiniteSquareCorralPhysics}
  \uses{def:physics:phyx-mini-0583:phyxminiproblems-problemphyxmini0583-lengthinpicometers, def:physics:phyx-mini-0583:phyxminiproblems-problemphyxmini0583-squarecorralexperiment, def:physics:phyx-mini-0583:phyxminiproblems-problemphyxmini0583-infinitesquarecorralmodeamplitudeperpicometer}
  The infinite-corral eigenstate law and Born's rule. The with-density equation uses two-dimensional Lebesgue volume on numerical picometre coordinates, so Complex.normSq has the reciprocal-square-picometre density role
\end{definition}
\begin{proof}
  This is the declared model definition of Satisfies Infinite Square Corral Physics.
\end{proof}
\begin{definition}[Matches Problem Statement And Figure]
  \label{def:physics:phyx-mini-0583:phyxminiproblems-problemphyxmini0583-matchesproblemstatementandfigure}
  \lean{PhyXMiniProblems.ProblemPhyXMini0583.MatchesProblemStatementAndFigure}
  \uses{def:physics:phyx-mini-0583:phyxminiproblems-problemphyxmini0583-lengthinpicometers, def:physics:phyx-mini-0583:phyxminiproblems-problemphyxmini0583-signedlengthinpicometers, def:physics:phyx-mini-0583:phyxminiproblems-problemphyxmini0583-squarecorralexperiment}
  Exact prose data and qualitative readouts from 583.png. These fields specify the apparatus and the E\_(1,3) state, but contain no detection probability, rounding claim, or answer-choice value
\end{definition}
\begin{proof}
  This is the declared model definition of Matches Problem Statement And Figure.
\end{proof}
\begin{definition}[detection Probability]
  \label{def:physics:phyx-mini-0583:phyxminiproblems-problemphyxmini0583-detectionprobability}
  \lean{PhyXMiniProblems.ProblemPhyXMini0583.detectionProbability}
  \uses{def:physics:phyx-mini-0583:phyxminiproblems-problemphyxmini0583-squarecorralexperiment, def:physics:phyx-mini-0583:phyxminiproblems-problemphyxmini0583-squareproberegionpicometers}
  Dimensionless Born probability that the electron lies in the probe region
\end{definition}
\begin{proof}
  This is the declared model definition of detection Probability.
\end{proof}
\begin{definition}[Answer Choice]
  \label{def:physics:phyx-mini-0583:phyxminiproblems-problemphyxmini0583-answerchoice}
  \lean{PhyXMiniProblems.ProblemPhyXMini0583.AnswerChoice}
  The four labels displayed beside the multiple-choice probabilities
\end{definition}
\begin{proof}
  This is the declared model definition of Answer Choice.
\end{proof}
\begin{definition}[displayed Probability]
  \label{def:physics:phyx-mini-0583:phyxminiproblems-problemphyxmini0583-answerchoice-displayedprobability}
  \lean{PhyXMiniProblems.ProblemPhyXMini0583.AnswerChoice.displayedProbability}
  \uses{def:physics:phyx-mini-0583:phyxminiproblems-problemphyxmini0583-answerchoice}
  Dimensionless probability printed for each answer choice
\end{definition}
\begin{proof}
  This is the declared model definition of displayed Probability.
\end{proof}
\begin{definition}[recorded Answer Choice]
  \label{def:physics:phyx-mini-0583:phyxminiproblems-problemphyxmini0583-recordedanswerchoice}
  \lean{PhyXMiniProblems.ProblemPhyXMini0583.recordedAnswerChoice}
  \uses{def:physics:phyx-mini-0583:phyxminiproblems-problemphyxmini0583-answerchoice}
  Dataset metadata records C; this name is never used as a premise
\end{definition}
\begin{proof}
  This is the declared model definition of recorded Answer Choice.
\end{proof}
\begin{definition}[Rounds To Nearest Ten Thousandth]
  \label{def:physics:phyx-mini-0583:phyxminiproblems-problemphyxmini0583-roundstonearesttenthousandth}
  \lean{PhyXMiniProblems.ProblemPhyXMini0583.RoundsToNearestTenThousandth}
  Rounding to the nearest 10\textasciicircum{}(-4), the precision of the displayed value 1.4 * 10\textasciicircum{}(-3). The lower endpoint is included and the upper excluded
\end{definition}
\begin{proof}
  This is the declared model definition of Rounds To Nearest Ten Thousandth.
\end{proof}
\begin{definition}[Matches Answer Choice]
  \label{def:physics:phyx-mini-0583:phyxminiproblems-problemphyxmini0583-matchesanswerchoice}
  \lean{PhyXMiniProblems.ProblemPhyXMini0583.MatchesAnswerChoice}
  \uses{def:physics:phyx-mini-0583:phyxminiproblems-problemphyxmini0583-squarecorralexperiment, def:physics:phyx-mini-0583:phyxminiproblems-problemphyxmini0583-detectionprobability, def:physics:phyx-mini-0583:phyxminiproblems-problemphyxmini0583-answerchoice, def:physics:phyx-mini-0583:phyxminiproblems-problemphyxmini0583-roundstonearesttenthousandth}
  A choice matches when the modeled probability rounds to its display
\end{definition}
\begin{proof}
  This is the declared model definition of Matches Answer Choice.
\end{proof}
\begin{lemma}[detection Probability rounds to one point four times ten neg three]
  \label{lem:physics:phyx-mini-0583:phyxminiproblems-problemphyxmini0583-detectionprobability-rounds-to-one-point-four-times-ten-neg-three}
  \lean{PhyXMiniProblems.ProblemPhyXMini0583.detectionProbability_rounds_to_one_point_four_times_ten_neg_three}
  \uses{def:physics:phyx-mini-0583:phyxminiproblems-problemphyxmini0583-squarecorralexperiment, def:physics:phyx-mini-0583:phyxminiproblems-problemphyxmini0583-satisfiesinfinitesquarecorralphysics, def:physics:phyx-mini-0583:phyxminiproblems-problemphyxmini0583-matchesproblemstatementandfigure, def:physics:phyx-mini-0583:phyxminiproblems-problemphyxmini0583-detectionprobability, def:physics:phyx-mini-0583:phyxminiproblems-problemphyxmini0583-roundstonearesttenthousandth}
  For the stated square corral, probe, and E\_(1,3) eigenstate, Born's rule gives a probability which rounds to 1.4 * 10\textasciicircum{}(-3)
\end{lemma}
\begin{proof}
  The conclusion follows from the governing relations and figure readouts named in the statement of detection Probability rounds to one point four times ten neg three.
\end{proof}
% --- Archon declaration topology end ---
% --- Archon physics formalization source end ---
